%==============================================================================
\section{BIRCH Clustering}
\label{sec:birch}
%==============================================================================

\begin{dinhnghia}[title={BIRCH}]
\textbf{B}alanced \textbf{I}terative \textbf{R}educing and \textbf{C}lustering \textbf{H}ierarchies:
thuật toán phân cụm phân cấp cho \textbf{dataset lớn}; xây cây cụm bằng chia đệ quy đến khi đạt tiêu chí dừng.
\end{dinhnghia}

\begin{ynghia}[title={CF và SCF}]
\begin{itemize}
  \item \textbf{Clustering Feature (CF)}: tóm tắt thống kê một tập điểm.
  \item \textbf{Sub-Cluster Feature (SCF)}: biểu diễn cấu trúc các subcluster.
\end{itemize}
\end{ynghia}

\subsection{Ba bước chính}

\begin{phuongphap}[title={Ba giai đoạn BIRCH}]
\begin{enumerate}
  \item \textbf{Khởi tạo}: tạo cây rỗng; đặt số CF tối đa mỗi nút.
  \item \textbf{Phân cụm}: đọc từng điểm; cập nhật CF nếu có, hoặc tạo CF mới; nếu vượt ngưỡng CF trong nút thì \textbf{chia subcluster} đệ quy.
  \item \textbf{Tinh chỉnh (Refinement)}: gộp subcluster tương tự theo metric khoảng cách.
\end{enumerate}
\end{phuongphap}

\subsection{Triển khai Python (scikit-learn)}

\begin{vidu}[title={Ví dụ Birch}]
\begin{lstlisting}
from sklearn.datasets import make_blobs
from sklearn.cluster import Birch
import matplotlib.pyplot as plt

X, y = make_blobs(n_samples=1000, centers=10, cluster_std=0.50,
                  random_state=0)

birch = Birch(threshold=1.5, n_clusters=4)
birch.fit(X)
labels = birch.predict(X)

plt.figure(figsize=(7.5, 3.5))
plt.scatter(X[:, 0], X[:, 1], c=labels, cmap='winter')
plt.show()
\end{lstlisting}
\end{vidu}

\begin{ynghia}[title={Tham số}]
\texttt{threshold}: ngưỡng tạo subcluster.
\texttt{n\_clusters}: số cụm đầu ra sau tinh chỉnh.
Dùng \texttt{fit} rồi \texttt{predict} để lấy nhãn.
\end{ynghia}

\tsimg{11_birch/img01.jpg}

\subsection{Ưu và nhược điểm}

\begin{tinhchat}[title={Ưu điểm}]
\begin{itemize}
  \item \textbf{Khả mở rộng} trên dữ liệu lớn nhờ cấu trúc cây.
  \item \textbf{Tiết kiệm bộ nhớ} nhờ CF/SCF.
  \item \textbf{Nhanh} nhờ phân cụm gia tăng (incremental).
\end{itemize}
\end{tinhchat}

\begin{luuy}[title={Nhược điểm}]
\begin{itemize}
  \item \textbf{Nhạy} tham số (số CF tối đa mỗi nút, ngưỡng subcluster).
  \item Giả định cụm \textbf{gần hình cầu}; kém với cụm phi cầu.
  \item Mặc định dùng \textbf{Euclidean}; có thể không phù hợp mọi bài toán.
\end{itemize}
\end{luuy}
